\documentclass[11pt,a4paper]{article}
\usepackage[utf8]{inputenc}
\usepackage[T1]{fontenc}
\usepackage{lmodern}
\usepackage[french]{babel}
\usepackage[landscape,margin=1.6cm,headheight=14pt]{geometry}
\usepackage{amsmath}
\usepackage{icomma}
\usepackage{booktabs,array,tabularx}
\usepackage{xcolor}
\usepackage{fancyhdr}
\usepackage{lastpage}
\usepackage[european,straightvoltages]{circuitikz}
\usetikzlibrary{calc,babel}

\definecolor{bleutitre}{HTML}{1F4E9A}
\pagestyle{fancy}
\fancyhf{}
\fancyhead[L]{\small Schéma électrique de la sonde de conductivité, tour hydroponique HydroTour}
\fancyhead[R]{\small Enoaarii RAPARII}
\fancyfoot[C]{\small Page \thepage\ / \pageref{LastPage}}
\renewcommand{\headrulewidth}{0.4pt}
\newcommand{\U}[1]{\,\mathrm{#1}}
\setlength{\parindent}{0pt}

\begin{document}

{\Large\bfseries Schéma électrique de la sonde de conductivité}

\vspace{1mm}
\begin{center}
\resizebox{0.88\linewidth}{!}{%
\begin{circuitikz}[line width=0.7pt, font=\small]
  % ---------------- Rails communs de l'Arduino ----------------
  \draw (0,7) node[vcc]{$+5\U{V}$} -- (19.6,7);
  \draw (0,0) node[ground]{} -- (19.6,0);
  \node[anchor=west, font=\footnotesize] at (19.7,7) {rail $+5\U{V}$ de l'Arduino};
  \node[anchor=west, font=\footnotesize] at (19.7,0) {masse (GND) de l'Arduino};

  % ---------------- Broche D5 ----------------
  \node[spdt, xscale=-1, anchor=in] (S1) at (3.5,4) {};
  \draw (S1.out 1) -- ++(-0.5,0) coordinate (a1) -- (a1 |- 0,7) node[circ]{};
  \draw (S1.out 2) -- ++(-0.5,0) coordinate (b1) -- (b1 |- 0,0) node[circ]{};
  \node[font=\bfseries, bleutitre] at (3.0,4.75) {D5};

  % ---------------- Pont diviseur et sonde ----------------
  \draw (S1.in) to[R, l={$R_{\text{ref}} = 330\U{\Omega}$}, i>^={$I$}] (7.5,4) node[circ] (J) {};
  \node[anchor=south east] at (7.4,4.05) {$J$};
  \draw (J) to[R, l={$R_w$}, v_>={$V_{w}$}] (11.5,4);
  \draw[dashed, teal!70!black] (8.3,2.75) rectangle (11.2,5.4);
  \node[teal!60!black, align=center, font=\footnotesize\bfseries] at (9.75,5.8) {Sonde graphite (eau du bac)};

  % ---------------- Broche D6 ----------------
  \node[spdt, yscale=-1, anchor=in] (S2) at (11.5,4) {};
  \draw (S2.out 2) -- ++(0.5,0) coordinate (a2) -- (a2 |- 0,7) node[circ]{};
  \draw (S2.out 1) -- ++(0.5,0) coordinate (b2) -- (b2 |- 0,0) node[circ]{};
  \node[font=\bfseries, bleutitre] at (12.0,3.15) {D6};

  % ---------------- Entrée A0 ----------------
  \draw (J) -- (7.5,2.2) node[ocirc] (A0) {};
  \node[anchor=west, font=\bfseries, bleutitre] at (7.65,2.2) {A0};
  \node[anchor=west, font=\footnotesize] at (7.65,1.75) {(CAN 10 bits)};
  \draw (7.5,0) node[circ]{} to[open, v<={$V_{A0}$}] (A0);

  % ---------------- DS18B20 ----------------
  \draw (15,7) node[circ]{} to[R, l_={$4{,}7\U{k\Omega}$}] (15,3.5) node[circ] (DQ) {};
  \draw (DQ) -- (14.1,3.5) node[ocirc]{};
  \node[anchor=east, font=\bfseries, bleutitre] at (14.0,3.5) {D8};
  \draw (DQ) -- (16.5,3.5);
  \draw[line width=1pt] (16.5,2.2) rectangle (18.9,4.6);
  \node[font=\footnotesize\bfseries] at (17.75,3.9) {DS18B20};
  \node[font=\scriptsize, anchor=west] at (16.55,3.5) {DQ};
  \draw (17,7) node[circ]{} -- (17,4.6);
  \node[font=\scriptsize, anchor=north] at (17,4.6) {VDD};
  \draw (18.2,2.2) -- (18.2,0) node[circ]{};
  \node[font=\scriptsize, anchor=south] at (18.2,2.2) {GND};
  \node[teal!60!black, font=\footnotesize\bfseries, anchor=east] at (18.05,1.6) {sonde dans l'eau};


  % ---------------- Repères ----------------
  \tikzset{rep/.style={circle, draw=black, fill=yellow!25, inner sep=1.2pt, minimum size=13pt, font=\scriptsize\bfseries}}
  \node[rep] at (1.3,7.4) {1};
  \node[rep] at (1.3,0.4) {2};
  \node[rep] at (2.35,4.75) {3};
  \node[rep] at (5.5,3.35) {4};
  \node[rep] at (7.9,4.55) {5};
  \node[rep] at (8.9,2.2) {6};
  \node[rep] at (8.65,5.05) {7};
  \node[rep] at (12.0,2.55) {8};
  \node[rep] at (15.6,5.25) {9};
  \node[rep] at (13.75,2.95) {10};
  \node[rep] at (19.3,4.3) {11};
  \node[font=\footnotesize\itshape] at (9.5,-1.3) {Inverseurs D5 et D6 représentés en phase 1 (mesure) : D5 relié au $+5\U{V}$, D6 relié à la masse.};
\end{circuitikz}}
\end{center}

\vspace{2mm}
{\small
\renewcommand{\arraystretch}{1.15}
\begin{tabularx}{\linewidth}{c >{\raggedright\arraybackslash}p{5.2cm} >{\raggedright\arraybackslash}X}
\toprule
\textbf{N°} & \textbf{Composant} & \textbf{Fonction} \\
\midrule
1 & Rail $+5\U{V}$ (broche 5V de l'Arduino) & Alimente la sonde EC quand D5 ou D6 est à l'état haut, alimente le DS18B20 et sa résistance de tirage. \\
2 & Masse (GND de l'Arduino) & Référence commune de toutes les tensions ; reçoit le courant quand D5 ou D6 est à l'état bas. \\
3 & Broche D5 (sortie numérique) & Applique $+5\U{V}$ puis $0\U{V}$ au pont diviseur ; en haute impédance au repos pour couper tout courant. \\
4 & $R_{\text{ref}} = 330\U{\Omega}$, $1\U{\%}$ & Résistance connue du pont diviseur ; limite le courant à $10\U{mA}$ au plus ; égale à $R_w$ au point de fonctionnement pour une sensibilité maximale. \\
5 & Point $J$ & Jonction entre $R_{\text{ref}}$ et la sonde : sa tension dépend de $R_w$. Soudé côté boîtier, jamais dans l'eau. \\
6 & Entrée A0 (convertisseur 10 bits) & Mesure $V_{A0}$ en phase 1 et la convertit en un entier $N$ de 0 à 1023 ; ne prélève pas de courant. \\
7 & Sonde graphite $R_w$ & Deux mines de graphite dans l'eau ; leur résistance $R_w = K / \sigma$ est la grandeur mesurée. \\
8 & Broche D6 (sortie numérique) & Applique $0\U{V}$ puis $+5\U{V}$, en opposition avec D5 : inverse le sens du courant (moyenne nulle, pas d'électrolyse). \\
9 & Résistance de tirage $4{,}7\U{k\Omega}$ & Maintient la ligne de données au $+5\U{V}$ au repos ; le DS18B20 et l'Arduino ne font que la tirer à la masse pour communiquer. \\
10 & Broche D8 (bus 1-Wire) & Dialogue numérique avec le DS18B20 : demande et lecture de la température. \\
11 & Capteur DS18B20 & Mesure la température de l'eau, pour ramener la conductivité à $25\,^{\circ}\mathrm{C}$. \\
\bottomrule
\end{tabularx}
}

\newpage
\vspace{3mm}
\begin{minipage}[t]{0.47\linewidth}
\textbf{Fonctionnement des broches}

\vspace{1mm}
\begin{tabular}{l c c c l}
\toprule
\textbf{Phase} & \textbf{D5} & \textbf{D6} & \textbf{$V_{D5} - V_{D6}$} & \textbf{Action} \\
\midrule
1 : mesure & $+5\U{V}$ & $0\U{V}$ & $+5\U{V}$ & lecture de $V_{A0}$ \\
2 : inversion & $0\U{V}$ & $+5\U{V}$ & $-5\U{V}$ & aucune lecture \\
Repos & \multicolumn{2}{c}{haute impédance} & aucune & aucun courant \\
\bottomrule
\end{tabular}

\vspace{2mm}
Chaque phase dure $212\U{\mu s}$ : signal carré de moyenne nulle à $2{,}36\U{kHz}$, répété 200 fois par mesure ($85\U{ms}$).
\end{minipage}
\hfill
\begin{minipage}[t]{0.49\linewidth}
\textbf{Lecture du schéma}

\vspace{1mm}
\begin{itemize}\setlength{\itemsep}{1pt}
\item Chaque broche numérique est représentée par un inverseur : elle relie le circuit soit au $+5\U{V}$, soit à la masse de l'Arduino.
\item $R_w$ est la résistance de l'eau entre les deux mines de graphite (environ $326\U{\Omega}$ à $1{,}5\U{mS \cdot cm^{-1}}$).
\item En phase 1, l'Arduino lit en A0 la tension du point $J$ par rapport à la masse : $V_{A0} = V_w$.
\item Résistance de l'eau : $R_w = R_{\text{ref}} \, N / (1024 - N)$.
\item Le point $J$ est soudé côté boîtier, jamais dans l'eau. Le DS18B20 reste sur D8 avec sa résistance de tirage de $4{,}7\U{k\Omega}$.
\end{itemize}
\end{minipage}

\end{document}
